% Bozza AI Eni 33545 - formato CV Barrington James 64a44d2, 29/10/2026.
% Contenuti Eni mantenuti; aggiornamento di formato richiesto il 02/11/2026.
\documentclass[10pt,a4paper]{article}


%=======================================================================%
%      1. PREAMBOLO - PACCHETTI E IMPAGINAZIONE
%=======================================================================%

\usepackage[utf8]{inputenc}
\usepackage[T1]{fontenc}
\usepackage[italian]{babel}
\usepackage{lmodern}
\usepackage[margin=1.7cm,headheight=14pt]{geometry}
\usepackage{xcolor}
\usepackage{enumitem}
\usepackage{tabularx}
\usepackage{microtype}
\usepackage{needspace}
\usepackage{fancyhdr}
\usepackage{lastpage}
\usepackage{fontawesome5}
\usepackage{hyperref}

\definecolor{accent}{HTML}{17365D}
\definecolor{muted}{HTML}{4A4A4A}

\hypersetup{
    colorlinks=true,
    urlcolor=accent,
    linkcolor=accent,
    pdfauthor={Antonio Michele Miti},
    pdftitle={Antonio Michele Miti - Eni HPC 33545 - Bozza}
}

\setlength{\parindent}{0pt}
\setlength{\parskip}{0pt}
\setlength{\emergencystretch}{2em}
\setlist[itemize]{
    leftmargin=1.3em,
    itemsep=2pt,
    topsep=3pt,
    parsep=0pt,
    partopsep=0pt
}

\pagestyle{fancy}
\fancyhf{}
\renewcommand{\headrulewidth}{0pt}
\cfoot{%
    \footnotesize\textcolor{muted}{%
        Antonio Michele Miti \enspace | \enspace
        Pagina \thepage\ di \pageref{LastPage}%
    }%
}


%=======================================================================%
%      2. MACRO - ELEMENTI RIUTILIZZABILI
%=======================================================================%

% Titolo di sezione con filetto e spazio minimo per il contenuto.
% Usage: \cvsection{Title}
\newcommand{\cvsection}[1]{%
    \Needspace{4\baselineskip}%
    \vspace{0.8em}%
    {\large\bfseries\color{accent} #1}\par
    \vspace{0.2em}\hrule\vspace{0.45em}%
}

% Voce con titolo, data, istituzione e descrizione.
% Usage: \entry{Title}{Date}{Institution}{Description}
\newcommand{\entry}[4]{%
    \Needspace{4\baselineskip}%
    \textbf{#1}\hfill #2\par
    \textit{#3}\par
    #4\par\vspace{0.45em}%
}

% Titolo di studio con voto sotto la data, nella colonna destra.
% Usage: \education{Title}{Date}{Grade}{Institution}{Description}
\newcommand{\education}[5]{%
    \Needspace{5\baselineskip}%
    \begin{tabularx}{\textwidth}{@{}>{\raggedright\arraybackslash}Xr@{}}
        \textbf{#1} & #2 \\
        \textit{#4} & #3
    \end{tabularx}\par
    #5\par\vspace{0.45em}%
}

% Progetto con rimando alle evidenze e lista dei contributi.
% Usage: \project{Title}{Evidence link}{Items}
\newcommand{\project}[3]{%
    \Needspace{5\baselineskip}%
    \textbf{#1}\par
    #2\par
    \begin{itemize}
        #3
    \end{itemize}
    \vspace{0.25em}%
}


%=======================================================================%
%      3. DOCUMENTO - PROFILO E PERCORSO PROFESSIONALE
%=======================================================================%

\definecolor{reviewred}{HTML}{B00020}
\newcommand{\AI}[1]{{\color{reviewred}#1}}
\begin{document}

    {\Huge\bfseries Antonio Michele Miti}\par
    \vspace{0.25em}
    {\large\color{accent} Fisica matematica e calcolo scientifico}\par
    \vspace{0.1em}
    \begin{center}
        \href{mailto:antoniomichele.miti@gmail.com}{%
                \faEnvelope\enspace antoniomichele.miti@gmail.com}  |
        \href{tel:+393339825418}{\faPhone\enspace +39 333 9825 418} \\[3pt]
        \href{https://www.antoniomiti.it/}{\faGlobe\enspace antoniomiti.it} |
        \href{https://github.com/MasterToninus}{\faGithub\enspace MasterToninus} |
        \href{https://orcid.org/0000-0002-8829-1943}{%
            \faOrcid\enspace ORCID: 0000-0002-8829-1943}
    \end{center}
    \vspace{0.1em}
    \textit{Candidatura: Junior Architetture Hardware di Supercalcolo-HPC}\par
    \textit{Eni -- San Donato Milanese -- posizione 33545}

{\footnotesize\color{reviewred} Bozza di revisione: rosso = nuova formulazione AI.}
\cvsection{Profilo}
\AI{Laureato magistrale in Fisica teorica e dottore di ricerca in Matematica, con un percorso internazionale nella fisica matematica. Affianco alla ricerca teorica esperienze universitarie di programmazione C++/CUDA e mantenimento di sistemi Linux e di un cluster didattico.}
\AI{Desidero applicare l’esperienza maturata nella fisica matematica a contesti industriali, attraverso un percorso di formazione continua e apprendimento sul campo delle tecnologie HPC.}
\cvsection{Competenze pertinenti}
\begin{tabularx}{\textwidth}{@{}>{\raggedright\arraybackslash\bfseries}p{3.5cm}>{\raggedright\arraybackslash}X@{}}
Calcolo scientifico & C++ e CUDA in progetti universitari: implementazione di algoritmi, confronto CPU--GPU, analisi delle prestazioni e stabilità numerica.\\[4pt]
Sistemi e strumenti & Linux, mantenimento di sistemi dipartimentali e cluster didattico; supporto agli utenti; Git e \LaTeX.\\[4pt]
Altri linguaggi & Python, Mathematica, Java e Haskell, utilizzati in contesti di studio, formazione e ricerca.\\[4pt]
Metodo e comunicazione & Formalizzazione dei problemi, verifica matematica, documentazione tecnica, didattica e collaborazione scientifica internazionale.
\end{tabularx}
\cvsection{Progetti di calcolo scientifico}
\entry{Calcolo di determinanti su CPU e GPU}{Progetto universitario}{C++ / CUDA}
{Implementazione ibrida di un algoritmo di condensazione per il calcolo del determinante. Confronto tra implementazione sequenziale e kernel GPU al variare della dimensione delle matrici; studio degli accessi alla memoria e del rapporto tra prestazioni e stabilità numerica.
\href{https://github.com/MasterToninus/Computational-Physics-Exercise}{Codice: Computational-Physics-Exercise}.}
\entry{Simulazione Monte Carlo del modello di Ising}{Progetto universitario}{C++}
{Implementazione e confronto degli algoritmi di Metropolis e a cluster mediante esperimenti numerici sullo stesso sistema fisico.}
\cvsection{Esperienze di supporto informatico}
\entry{Collaborazione studentesca -- mantenimento cluster}{05/2014--02/2015}{Laboratorio di Calcolo e Multimedia, Università degli Studi di Milano}
{Mantenimento del cluster del laboratorio e supporto agli utenti e alle attività didattiche di fisica computazionale in ambiente Linux.}
\entry{Supporto ai sistemi informatici dipartimentali}{03/2017--10/2018}{Università Cattolica del Sacro Cuore, Brescia}
{Attività di assistenza al mantenimento dei sistemi informatici del Dipartimento di Matematica e Fisica, durante il dottorato.}
\entry{Stage di programmazione Java}{03/2016--07/2016}{CST s.r.l., Milano}
{Formazione in Java, programmazione orientata agli oggetti e strumenti per lo sviluppo software professionale.}
\newpage
\cvsection{Esperienza nella ricerca}
\entry{Postdoctoral Fellow -- CIVIS3i, MSCA COFUND}{2024--2026}{Sapienza Università di Roma}
{Progetto MultiSymORI su geometria multisimplettica, osservabili e riduzione, nell'ambito della fisica matematica. Ricerca, collaborazioni internazionali, visite scientifiche, pubblicazioni e rendicontazione.}
\entry{Assegnista di ricerca}{10/2022--01/2024}{Università Cattolica del Sacro Cuore, Brescia}
{Ricerca in geometria simplettica e multisimplettica con applicazioni alla fisica matematica; collaborazione internazionale e didattica universitaria.}
\entry{Research Fellow}{10/2021--09/2022}{Max Planck Institute for Mathematics, Bonn}
{Ricerca sulle strutture algebriche superiori in geometria e fisica matematica, in un ambiente internazionale.}
\entry{Dottorando in Matematica}{2016--2021}{Università Cattolica del Sacro Cuore e KU Leuven}
{Dottorato internazionale con periodi di ricerca in Italia e Belgio; geometria differenziale, meccanica geometrica e teorie classiche dei campi.}
\cvsection{Formazione}
\education{Dottorato in Science: Mathematics}{04/2021}{}{Università Cattolica del Sacro Cuore / KU Leuven}
{Titolo congiunto internazionale. Tesi: \emph{Homotopy comomentum maps in multisymplectic geometry}.}
\education{Laurea magistrale in Fisica -- curriculum teorico}{11/2015}{110/110}{Università degli Studi di Milano}
{Metodi matematici della fisica, meccanica geometrica e metodi computazionali.}
\education{Laurea in Fisica}{10/2010}{99/110}{Università degli Studi di Milano-Bicocca}{}
\cvsection{Transizione professionale e formazione tecnica}
\AI{Iscritto alla selezione del Master universitario di II livello in \emph{High-Performance and Quantum Computing}, Università di Bologna, a.a.\ 2026/2027 (codice 6162).}\par
\begin{tabularx}{\textwidth}{@{}p{1.2cm}X@{}}
2024 & Workshop \emph{High Performance Computing and Quantum Computing}, VII edizione, CINECA.\\[3pt]
2019 & \emph{Functional Programming and Domain-Specific Languages}, KU Leuven (Haskell).\\[3pt]
2018 & \emph{Scientific Computing with Python}, Università Cattolica.\\[3pt]
2017 & \emph{Python for Computational Science}, CINECA-SCAI.\\[3pt]
2016 & \emph{Introduction to Scientific and Technical Computing in C++}, CINECA-SCAI.\\[3pt]
2014 & \emph{Scientific Programming using GPGPU and CUDA}, CINECA-SCAI.
\end{tabularx}
\cvsection{Pubblicazione selezionata}
A. Luongo, A. M. Miti, V. Narasimhachar, A. Sireesh.
\emph{Measurement-based uncomputation of quantum circuits for modular arithmetic}. DAC 2025.
\href{https://doi.org/10.1109/DAC63849.2025.11132981}{DOI}.\par
Contributo congiunto a costruzioni di circuiti aritmetici quantistici e alla loro verifica matematica; confronto dei costi in risorse. Esperienza complementare al calcolo classico CPU--GPU.
\cvsection{Lingue e disponibilità}
Italiano: madrelingua. Inglese: uso professionale nella ricerca e nella didattica; TOEFL iBT 102/120 (aprile 2018, risultato storico). Francese: base.\par
Disponibile a lavorare a San Donato Milanese e alla mobilità nazionale e internazionale.
\end{document}
